\section{Точная огибающая полного зазора}
\label{boundarygap:section}
При $g>0$ определим
\begin{equation}\label{boundarygap:definition}
 V(g,\rho)=A(g/\sqrt\rho,\sqrt\rho)\quad(0<\rho\le g/m),
 \qquad V(g,0)=g,\qquad
 \mathcal B(g)=\max_{0\le\rho\le g/m}V(g,\rho),
\end{equation}
и положим $\mathcal B(0)=0$.
Здесь $\rho=0$ обозначает предел сжимающихся областей, не точку
с положительной ошибкой на нулевом расстоянии от минимума.
Другой конец даёт $V(g,g/m)=W(g)=\mathcal E(\sqrt{mg})$.
Равносильная радиальная запись:
\begin{equation}\label{joint:radial-identity}
 A(P,r)=V(Pr,r^2),\qquad A(P,P/m)=\mathcal E(P).
\end{equation}

\begin{theorem}[Сертификат полного зазора]
\label{boundarygap:main}
При условиях \eqref{boundarygap:assumptions}
\[
 \delta(x)\le\mathcal B(G(x)).
\]
Максимум в \eqref{boundarygap:definition} существует.
Функция $\mathcal B$ конечна, непрерывна и удовлетворяет
\begin{equation}\label{boundarygap:growth}
 \mathcal B(g_2)\ge\mathcal B(g_1)+g_2-g_1
 \quad(0\le g_1<g_2),\qquad \mathcal B(g)\ge g.
\end{equation}
Следовательно, она строго возрастает, не ограничена сверху
и имеет непрерывную обратную функцию.
При фиксированном точном зазоре $G(x)=g>0$ значение
$\mathcal B(g)$ есть супремум ошибки. Его приближают гладкие
функции на отрезках с теми же константами и зазором.

При $H_1=0$, с $\beta$ из \eqref{compact:radial-constant},
\begin{equation}\label{boundarygap:constant}
 V(g,\rho)=g+\beta\rho,\qquad
 \mathcal B(g)=g\max\{1,1+\beta/m\}.
\end{equation}

При $H_1>0$ положим $T_g=\sqrt{H_1g/m}$ и
\begin{equation}\label{boundarygap:elementary-formula}
 \mathcal B(g)\le\mathcal C(g):=
 \frac{H_0}{H_1}(\cosh T_g-1)+g\frac{\sinh T_g}{T_g},
 \qquad\mathcal C(0)=0.
\end{equation}
Функция $\mathcal C$ непрерывна, строго возрастает и не ограничена
сверху; её вычисление не требует корня или внешней максимизации.
\end{theorem}

\begin{proof}\label{boundarygap:proof}
Для $g=G(x)>0$ радиальные данные дают $r^2\le g/m$ и
$u'(r)\le g/r$. По лемме~\ref{compact:radial}
$\delta(x)\le A(g/r,r)=V(g,r^2)$.
При $g=0$ точка оптимальна.
Непрерывность при $\rho>0$ следует из уравнения переключения.
Для $H_1>0$, $c=\sqrt{H_1\rho}$, интегрирование RSI и
сравнение с положительной кривизной на всём пути дают
\begin{equation}\label{compact:radial-upper}
 g-\frac{m\rho}{2}\le V(g,\rho)
 \le\frac{H_0}{H_1}(\cosh c-1)
        +(g-m\rho)\frac{\sinh c}{c}.
\end{equation}
Обе границы стремятся к $g$ при $\rho\downarrow0$ локально
равномерно по $g>0$. После замены $\rho=zg/m$, $z\in[0,1]$,
максимум берётся по фиксированному компакту и непрерывен по $g$.
Правая часть имеет порядок $O(g)$ равномерно по $z$ при $g\downarrow0$.
Случай $H_1=0$ следует из явной формулы.
Верхняя оценка $\mathcal C$ получается заменой $g-m\rho$ на $g$
и $\rho$ на $g/m$ в \eqref{compact:radial-upper}.
Остальные её свойства следуют из степенных рядов с неотрицательными
коэффициентами и $\mathcal C(g)\ge g$.

При положительном максимизирующем радиусе точность даёт
гладкое приближение из радиальной леммы: на $[0,r]$ полный
зазор в правом конце равен $r(g/r)=g$ точно.
Если максимум задаёт $\rho=0$, то функции
\[
 u(s)=(g/r-mr)s+ms^2/2,\qquad 0\le s\le r,
\]
имеют зазор $g$ и потери $g-mr^2/2\to g$.
Наконец, добавление $(g_2-g_1)s/r$ к почти экстремальной
функции для $g_1>0$ сохраняет сдвинутое RSI и гессиан,
увеличивает допустимую кривизну и повышает зазор и ошибку
ровно на $g_2-g_1$. Это доказывает \eqref{boundarygap:growth};
случай $g_1=0$ следует из сжимающихся примеров.
\end{proof}

\paragraph{Переключение в нормированном времени.}
Для вычислений при $H_1>0$, $g>0$, $0<\rho\le g/m$ положим
$c=\sqrt{H_1\rho}$, $a=H_0/H_1$, $q=g-m\rho$ и
\[
 \begin{aligned}
 Y_\tau&=a\cosh(c\tau)+(q/c)\sinh(c\tau),&
 Q_\tau&=ac\sinh(c\tau)+q\cosh(c\tau),\\
 F(\tau)&=Q_\tau\cos b-cY_\tau\sin b,&
 Z(\tau)&=Y_\tau\cos b+(Q_\tau/c)\sin b-a,
 \end{aligned}
 \qquad b=c(1-\tau).
\]
Радиальная лемма после замены времени даёт единственный внутренний
корень и значение
\begin{equation}\label{boundarygap:switch}
 F(\tau)=g,\qquad
 \max\{0,1-\pi/(4c)\}\le\tau\le1,\qquad V(g,\rho)=Z(\tau).
\end{equation}
Прямое дифференцирование даёт
\begin{equation}\label{boundarygap:switch-derivatives}
 F'=2\rho(H_0+H_1u(\tau))\cos b>0,\qquad
 Z'=\frac{2\rho(H_0+H_1u(\tau))}{c}\sin b\ge0,
\end{equation}
где $u(\tau)=Y_\tau-a$. Верхний конец интервала для $\tau$
поэтому даёт верхнюю оценку $V$ без численного интегрирования.

\begin{proposition}[Единый бюджет скалярного ограждения]
\label{boundarygap:budget}\label{joint:budget}
Пусть $H_1>0$, $g>0$, $\vartheta>1$. Определим
\[
 T=\sqrt{H_1g/m},\quad
 \mathsf a(t)=\frac{\sinh t}{t},\quad
 \mathsf c(t)=\frac{\cosh t-1}{t^2},\qquad
 \mathsf a(0)=1,\quad\mathsf c(0)=\tfrac12,
\]
\begin{equation}\label{compact:scalar-constants}
\begin{aligned}
 K_*&=H_0+H_1\mathcal C(g),\\
 L_*&=m\{\mathsf a(T)+2\cosh T\}
       +K_*\{\mathsf c(T)+2\mathsf a(T)\}.
\end{aligned}
\end{equation}
На отрезке $I=[\rho_-,\rho_+]\subseteq[0,g/m]$ пусть
$M=\max_I V(g,\rho)$ и известна граница $0<\ell\le M$.
Можно всегда взять $\ell=g/2$; при $I=[0,g/m]$ можно взять $\ell=g$.
Это нижняя граница огибающей $M$, а не ошибки конкретной точки.
Для $\Delta=\rho_+-\rho_->0$ положим
\begin{equation}\label{compact:interval-budget}
 J=\max\left\{1,\left\lceil\frac{L_*\Delta}
                       {(\vartheta-1)\ell}\right\rceil\right\},
 \qquad n=\left\lceil\log_2\max\left\{1,
       \frac{4\rho_+K_*}{(\vartheta-1)\ell}\right\}\right\rceil.
\end{equation}
В каждом положительном узле из $J+1$ равномерных узлов выполним $n$
делений пополам в \eqref{boundarygap:switch} с точным знаком.
Пусть $\widehat V_j$ есть $Z$ на верхнем сохранённом конце;
для узла $\rho_j=0$ положим $\widehat V_j=g$.
Тогда
\begin{equation}\label{boundarygap:finite-enclosure}
 \widehat M=\max_j\widehat V_j+\frac{L_*\Delta}{2J},
 \qquad M\le\widehat M\le\vartheta M.
\end{equation}
Положительных узлов $N_+=J+1-\mathbf1_{\{\rho_-=0\}}$;
нужно $N_+n$ сравнений знака и $J$ сравнений для максимума.
При $\Delta=0$, $\rho_+>0$ достаточно одного узла с
\begin{equation}\label{energy:scalar-budget}
 n=\left\lceil\log_2\max\left\{1,
       \frac{2\rho_+K_*}{(\vartheta-1)\ell}\right\}\right\rceil
\end{equation}
сравнениями знака, без сеточной поправки. Для $I=\{0\}$ ответ равен $g$.
\end{proposition}

\begin{proof}
Оценка \eqref{compact:radial-upper} даёт $V\ge g/2$.
При оптимальном переключении $0\le u\le\mathcal C(g)$.
Если сначала фиксировать переключение, то чувствительность
$w=\partial_\rho u$ решает
\[
 w''=\sigma_\tau(H_0+H_1u+\rho H_1w),\qquad
 w(0)=0,\quad w'(0)=-m.
\]
Сравнение интегрального уравнения по модулю с положительным ядром даёт
\[
 |w(1)|\le m\mathsf a(c)+K_*\mathsf c(c),\qquad
 |w'(1)|\le m\cosh c+K_*\mathsf a(c),\qquad c=\sqrt{H_1\rho}.
\]
После учёта условия постоянной конечной производной
\begin{equation}\label{boundarygap:rho-derivative}
 V_\rho=w(1)-(\tan b/c)w'(1).
\end{equation}
Поскольку $b\le\min\{c,\pi/4\}$, имеем $\tan b\le2c$,
откуда $|V_\rho|\le L_*$.
Непрерывность продолжает эту оценку на концы отрезка.
В \eqref{boundarygap:switch-derivatives} имеем
$0\le Z'\le2\rho K_*$, в том числе при пробных переключениях.
Поэтому ошибка после $n$ делений не больше
$2\rho_+K_*2^{-n}\le(\vartheta-1)\ell/2$.
Любая точка $I$ удалена от узла не более чем на $\Delta/(2J)$;
поправка не больше той же величины.
Их сумма не превосходит $(\vartheta-1)M$.
При одном узле весь допуск отводится ошибке корня.
\end{proof}

Это счёт сравнений в точной арифметике, не число интервальных
операций и не битовая сложность. Строгая реализация должна
ограждать значения элементарных функций и ошибки округления;
например, можно строить оценку с $1<\vartheta_0<\vartheta$
и оставить запас на эти ошибки. Новых запросов функции,
градиента или LMO скалярная процедура не делает.
При $H_1=0$ используются явные формулы, при $g=0$ возвращается нуль.
Если интервальный знак не определён, нужна другая корректная
процедура либо явный отказ в сертификате.
При $g\downarrow0$ имеем $K_*\to H_0$ и
$L_*\to3m+5H_0/2$, поэтому бюджеты для полного зазора и
совместного сертификата остаются ограниченными при фиксированных
константах и множителе.

\paragraph{Двух крайних значений недостаточно.}
\label{boundarygap:counterexample}
При $m=H_0=H_1=1$, $g=6$ формула
\eqref{gapregime:variation} даёт
\begin{equation}\label{boundarygap:small-radius}
 V(6,\rho)=6+\frac{181-31\sqrt{31}}{54}\rho+O(\rho^2)>6
\end{equation}
при малых положительных $\rho$.
При этом в стационарном интеграле времени для $U=6$, $P=\sqrt6$
пересечение равно $v_c=2\sqrt7-1$, а наибольшая скорость равна
$3\sqrt3$. Значит,
\[
 \mathcal T(6;\sqrt6,0)
 \ge\operatorname{arcosh}(2\sqrt7)
       +\frac{7-2\sqrt7}{3\sqrt3}
 >\frac{13}{6}+\frac3{10}>\sqrt6.
\]
Здесь $\operatorname{arcosh}(2\sqrt7)>\log9>13/6$
следует из первых двух членов ряда для $\log3$;
$\sqrt7<8/3$, $\sqrt3<7/4$ дают вторую строгую оценку.
По строгому обращению $W(6)<6$.
Следовательно, $\mathcal B(6)>\max\{6,W(6)\}$.
